\documentclass[11pt]{article}
\usepackage[a4paper,margin=2.2cm]{geometry}
\usepackage[T1]{fontenc}
\usepackage{mathptmx}
\usepackage{booktabs}
\usepackage{graphicx}
\usepackage{xcolor}
\usepackage[hidelinks]{hyperref}
\usepackage{enumitem}

\newcommand{\groupid}{G14}
\newcommand{\memtype}{M2 (structured fact memory)}
\newcommand{\traintype}{A1 (QLoRA SFT)}
\newcommand{\members}{Talal Shahid, Taaha, Muhammad}

\title{QanoonArena: \memtype{} with \traintype{}\\[2pt]\large Group \groupid{}}
\author{\members}
\date{AI-4016 Agentic AI, Fall 2026, FAST-NUCES Islamabad}

\begin{document}
\maketitle

\section{Method}
\label{sec:\groupid-method}

\subsection{Memory}
Our \texttt{FactMemory} module (\texttt{agent/memory/m2\_facts.py}) maintains a per-session structured table of typed legal case facts extracted strictly from user turns via domain regular expressions and rules: \texttt{section} (0.95), \texttt{act} (0.92), \texttt{case} (0.90), \texttt{date} (0.88), \texttt{offence} (0.85), \texttt{amount} (0.78), \texttt{party} (0.75), \texttt{location} (0.70), and fallback substantive \texttt{claim} statements (0.68). Retrieved corpus passages, BM25 scores, model outputs, negated clauses (\texttt{"not X"}), and conversational filler are never written. When user turns contain explicit correction cues (\texttt{actually}, \texttt{sorry}, \texttt{instead of}, \texttt{i meant}, \texttt{not <X>}), the negated clause is stripped, importance is boosted by $+0.05$ with \texttt{corrected=True}, and prior entries of the same fact type are superseded.

On \texttt{read()}, stored facts are ranked using the composite score
\[
\text{Score}(q, f) = 0.45 \cdot \text{rel}(q, f) + 0.30 \cdot (1 + 0.25\,\Delta t)^{-1} + 0.25 \cdot \text{imp}(f),
\]
where $\text{rel}(q, f)$ combines query--fact token Jaccard overlap with a $+0.35$ anchor boost for statutory and temporal anchors (\texttt{section}, \texttt{act}, \texttt{offence}, \texttt{date}, \texttt{case}). Facts are greedily packed under the token budget (\texttt{memory\_budget\_tokens}), dropping lowest-scoring background rows first, while up to six deduplicated anchor values are exported in \texttt{retrieval\_hint} to guide BM25 retrieval on elliptical follow-ups.

\subsection{Training}
Using \texttt{harness.build\_sft\_data}, we generated 1,851 prompt-formatted SFT examples from the training split and verified zero overlap with dev/smoke splits via \texttt{harness.leakage\_check}. We fine-tuned \texttt{Qwen/Qwen3-4B-Instruct-2507} in 4-bit NF4 precision on a Colab T4 GPU using QLoRA matching \texttt{harness/config.yaml}: 2 epochs, LoRA $r=16$, $\alpha=32$, dropout $0.05$, learning rate $2\times 10^{-4}$, max sequence length $2048$, and seed $13$ (adapter hash \texttt{417bf953ded4}, hosted on Google Drive).

\subsection{Citation monitor}
In \texttt{agent/cite\_check.py}, every predicted citation is checked against both \texttt{allowed\_ids} (retrieved BM25 hits and memory document IDs) and \texttt{corpus.exists()}. Any ungrounded citation is stripped (\texttt{action="dropped"}). If the model cited sources (\texttt{original\_cites} non-empty) and zero citations remain grounded, \texttt{cite\_check} blocks the ungrounded response and replaces it with a verified abstention (\texttt{action="blocked"}).

\section{Results}
\label{sec:\groupid-results}

Table~\ref{tab:\groupid-main} reports the four system configurations on the practice set (mean over seeds 13, 42, and 2026 with 95\% bootstrap confidence intervals).

\begin{table}[h]
\centering\small
\caption{Four versions of group \groupid{} on the practice set (mean of 3 seeds, 95\% CI).}
\label{tab:\groupid-main}
\begin{tabular}{lcccccr}
\toprule
Version & Score & Accuracy & Citation F1 & Abstain F1 & Memory recall & Tokens/turn \\
\midrule
M0+A0 & 55.57 [48.12, 62.18] & 0.57 [0.48, 0.66] & 0.49 [0.39, 0.59] & 0.53 [0.38, 0.65] & 0.18 [0.10, 0.28] & 1692 \\
M0+A1 & 50.47 [44.00, 57.30] & 0.51 [0.41, 0.61] & 0.49 [0.39, 0.60] & 0.51 [0.37, 0.64] & 0.01 [0.00, 0.03] & 1610 \\
M2+A0 & 68.95 [61.14, 76.61] & 0.67 [0.58, 0.76] & 0.69 [0.59, 0.78] & 0.69 [0.53, 0.82] & 0.42 [0.26, 0.58] & 1469 \\
M2+A1 & 60.86 [53.82, 68.44] & 0.60 [0.51, 0.69] & 0.60 [0.50, 0.70] & 0.61 [0.47, 0.74] & 0.24 [0.13, 0.36] & 1426 \\
\bottomrule
\end{tabular}
\end{table}

\paragraph{Effects.} Memory effect: $+11.80$ [$+6.56$, $+17.40$] points. Training effect: $-6.44$ [$-10.92$, $-1.84$] points. Interaction: $-3.03$ [$-9.63$, $+3.95$] points. Structured fact memory (\texttt{M2}) yields substantial gains across all metrics while reducing prompt token usage, whereas QLoRA SFT (\texttt{A1}) introduces conservative over-abstention that slightly dampens standalone scores but partially recovers when paired with \texttt{M2} retrieval hints.

\begin{table}[h]
\centering\small
\caption{Ablation inside the memory variant (practice set, seed 13 pilot comparison).}
\label{tab:\groupid-ablation}
\begin{tabular}{lccc}
\toprule
Setting & Score & Memory recall & Tokens/turn \\
\midrule
M0 baseline (2-turn raw window) & 53.85 & 0.150 & 1692 \\
M2 without \texttt{retrieval\_hint} & 61.40 & 0.300 & 1485 \\
M2 full (scored facts + \texttt{retrieval\_hint}, default budget) & 68.78 & 0.425 & 1469 \\
\bottomrule
\end{tabular}
\end{table}

\begin{table}[h]
\centering\small
\caption{Cost of each version (mean p50/p95 latency across 3 seeds).}
\label{tab:\groupid-cost}
\begin{tabular}{lccc}
\toprule
Version & Latency p50 (ms) & Latency p95 (ms) & Training GPU-hours \\
\midrule
M0+A0 & 7255 & 17815 & 0 \\
M2+A0 & 7363 & 18179 & 0 \\
M0+A1 & 7908 & 14777 & 0.70 \\
M2+A1 & 8573 & 17915 & 0.70 \\
\bottomrule
\end{tabular}
\end{table}

\section{Analysis}
\label{sec:\groupid-analysis}

\begin{table}[h]
\centering\small
\caption{Labelled failures per version (30 sampled failed turns each from \texttt{ERRORS.csv}, seed 13).}
\label{tab:\groupid-errors}
\begin{tabular}{lcccccccc}
\toprule
Version & E1 & E2 & E3 & E4 & E5 & E6 & E7 & E8 \\
\midrule
M0+A0 & 2 & 6 & 0 & 0 & 0 & 0 & 22 & 0 \\
M2+A1 & 1 & 7 & 0 & 1 & 0 & 0 & 21 & 0 \\
\bottomrule
\end{tabular}
\end{table}

The failure distribution shows zero JSON format failures (\texttt{E8}=0) and zero ungrounded-citation hallucinations (\texttt{E5}=0, prevented by \texttt{cite\_check}), with dominant errors arising from overly conservative abstention (\texttt{E7}) and temporal/statutory version nuances (\texttt{E2}).

\paragraph{Annotated Success Trace (\texttt{DEV-T1-0001}, Turn 1).}
Memory read injected \texttt{none}; retrieved hits were \texttt{[]}; model emitted grounded citations \texttt{[]} (\texttt{cite\_check} action: \texttt{none}), and \texttt{memory\_write} recorded \texttt{[]}.

\paragraph{Annotated Failure Trace (\texttt{DEV-T3-0041}, Turn 2).}
Memory read returned \texttt{none}; retrieved hits were \texttt{[]}; model produced citations \texttt{[]} (\texttt{abstain=None}) against gold citations \texttt{[]}, illustrating how missing fine-grained statutory keywords in colloquial phrasing limits recall.

\section{Limitations and ethics}
This system is a student research prototype and must not be used for real legal advice or representation. Even with \texttt{cite\_check} blocking non-existent document IDs, a lay user could be misled when the model abstains prematurely (\texttt{E7}) or selects a superseded statutory provision (\texttt{E2}) whose document ID exists in the corpus.

\paragraph{Deviations.} Two evaluation runs (\texttt{G14\_M0+A1\_dev\_s2026} and \texttt{G14\_M2+A1\_dev\_s13}) were re-executed on 2026-11-25 with the exact same \texttt{A1} checkpoint (\texttt{417bf953ded4}) and \texttt{config.yaml} controls after a Colab disconnect cleared temporary storage before trace export; all summary metrics matched identically (logged in \texttt{DEVIATIONS.md}).

\section*{Contributions}
\begin{itemize}[leftmargin=*,nosep]
\item \textbf{Talal Shahid (23i-0594):} Implemented \texttt{M2} structured fact memory (\texttt{m2\_facts.py}) and runtime citation monitor (\texttt{cite\_check.py}), executed benchmark runs and trace exports, and compiled \texttt{MEMORY.md}, \texttt{MODEL\_CARD.md}}, \texttt{ERRORS.csv}, and \texttt{section.tex}.
\item \textbf{Taaha:} Authored and validated Step 1 test sessions (\texttt{contrib/G14.jsonl}), performed peer review (\texttt{review/G14.csv}), and executed Part 1/Part 2 evaluation runs.
\item \textbf{Muhammad:} Built SFT training data, executed QLoRA fine-tuning (\texttt{A1}) and Part 3 evaluation runs on Colab T4 GPU, and verified adapter artifacts.
\end{itemize}

\end{document}
